\documentclass{article}
\usepackage[utf8]{inputenc}
\usepackage[T1]{fontenc}
\usepackage[english, polish]{babel}
\usepackage{amsmath}
\usepackage{amsfonts}
\usepackage{amssymb}
\usepackage{graphicx}
\usepackage{float}
\usepackage{parskip}
\usepackage{hyperref}
% Hide the link borders
\hypersetup{hidelinks}
\usepackage{geometry}
% Margins for A4 paper
\geometry{a4paper, margin=2cm}
\usepackage{enumitem}
% Spacing for itemize lists
\setlist[itemize]{itemsep=1em, topsep=0.25em, parsep=0pt, partopsep=0pt}
\setlength{\parindent}{0pt} % No paragraph indentation
\setlength{\emergencystretch}{2em} % Prevent overfull hboxes

\title{Sprawozdanie z zadania 1\\\textbf{Gra wyścigowa}}
\author{Mikołaj Kuryło s193743\\Antonina Włodarczyk s197793}
\date{}

\begin{document}
\maketitle
\vspace{1cm}

\section{Sformułowanie zadania i przyjęte założenia}
Celem projektu było stworzenie dwuwymiarowej gry wyścigowej z widokiem prostopadłym do płaszczyzny toru. Gra miała kłaść nacisk na realistyczny (uproszczony fizycznie) model jazdy samochodem.\par
Do założeń szczegółowych należało:
\begin{itemize}
    \item \textbf{Algorytm gry:} Gra opiera się na pętli sterowanej przez systemowy zegar (prywatna metoda \texttt{\_tick} w klasie \texttt{Game}), która jest wywoływana z częstotliwością 60 razy na sekundę lub mniej w zależności od wydajności komputera. W każdej iteracji pętli dochodzi do aktualizacji stanu gry, w tym fizyki ruchomych i stacjonarnych obiektów, a następnie do renderowania aktualnego stanu gry na ekranie.
    \item \textbf{Model fizyczny:} Samochód posiada masę, rozstaw osi oraz parametry takie jak przyczepność opon zależna od podłoża. W modelu uwzględniono siły odśrodkowe, opory powierza oraz tarcie, transfer masy między osiami będący wynikową przyspieszenia i hamowania.
    \item \textbf{Grafika:} Wszystkie graficzne elementy gry (samochody, kafelki terenu) są ładowane z plików graficznych (.png) i dynamicznie obracane w zależności od kierunku poruszania się pojazdu. Kafle terenu współdzielą mechanizm renderowania z samochodem z wyjątkiem tego, że się nie przesuwają, a jedynie są renderowane w odpowiednich miejscach na mapie. Mapa budowana jest z siatki kafelków (droga, trawa, piasek), z których każdy posiada własny współczynnik tarcia wpływający na zachowanie auta.
    \item \textbf{Sterowanie:} Gracz steruje samochodem za pomocą klawiatury, gdzie klawisze strzałek służą do przyspieszania, hamowania i skręcania. Sterowanie jest responsywne, a reakcja samochodu na polecenia gracza jest natychmiastowa.
\end{itemize}

\section{Opis rozwiązań programowych}
Aplikacja została napisana w języku \textit{Python} z wykorzystaniem biblioteki \textit{Tkinter} do obsługi grafiki (w tym renderowania i GUI), \textit{PIL} (Pillow) do manipulacji obrazami oraz \textit{NumPy} do obliczeń matematycznych związanych z fizyką ruchu samochodu.\par
\subsection{Architektura systemu}
Projekt został podzielony na moduły zapewniające separację logiki od prezentacji i ułatwiające zarządzanie kodem:
\begin{itemize}
    \item \texttt{logic.py} - rdzeń logiki gry, zawierający klasy i funkcje odpowiedzialne za fizykę, sterowanie, zarządzanie stanem gry, wczytywanie mapy i interakcje między obiektami. Implementuje dynamiczne sprawdzanie podłoża. W każdej klatce program rzutuje pozycję x, y samochodu na macierz \texttt{traction\_map\_matrix}. Dzięki temu współczynnik tarcia µ zmienia się płynnie między asfaltem a trawą, co bezpośrednio wpływa na sterowność.
    \item \texttt{renderer.py} - ze względu na ograniczenia wydajnościowe biblioteki \textit{Tkinter}, zastosowano autorski system cache'owania tekstur w klasie \texttt{Texture}. Zamiast obracać obrazek przy każdym wywołaniu pętli, obrócone instancje \texttt{PhotoImage} są przechowywane w słowniku, gdzie kluczem jest kąt (0-359 stopni). Redukuje to obciążenie CPU o rząd wielkości podczas intensywnych manewrów.
    \item \texttt{gui.py} i \texttt{menu.py} - moduły obsługujące interfejs użytkownika, okna startowe oraz wybór tekstur.
    \item \texttt{map\_reader.py} - moduł do wczytywania mapy z pliku .json, który zawiera informacje o układzie kafelków terenu oraz pozycjach startowych samochodów.
\end{itemize}
\subsection{Model fizyczny - fragment kodu}
Poniższy fragment kodu z metody \texttt{update} klasy \texttt{CarEntity} rozszerzający podstawową klasę \texttt{Entity} pokazuje implementację modelu fizycznego samochodu, w tym obliczanie sił działających na pojazd oraz aktualizację jego pozycji i prędkości:
\begin{verbatim}
    # --- Lateral forces calculations ---
    # Calculate Slip Angles
    # Front slip = atan(lat_vel + angular_vel * dist_to_front / long_vel) - steer_angle
    if abs(self.velocity.y) > const.MIN_VELOCITY_THRESHOLD:  # Avoid division by zero
        slip_angle_front = np.degrees(np.arctan2(
            self.velocity.x + self.angular_velocity * self.wheelbase *
            (1-self.mass_distribution), abs(self.velocity.y))) - self.steer_angle
        slip_angle_rear = np.degrees(np.arctan2(
            self.velocity.x - self.angular_velocity * self.wheelbase *
            self.mass_distribution, abs(self.velocity.y)))
    else:
        slip_angle_front = slip_angle_rear = 0.0

    # Calculate Lateral Forces
    # F = -stiffness * slip_angle
    force_front_lat = -self.cornering_stiffness * slip_angle_front
    force_rear_lat = -self.cornering_stiffness * slip_angle_rear

    # --- Longitudinal forces ---
    force_front_long = accel_force * front_accel_force_long_coeff - brake_force
    force_rear_long = accel_force * rear_accel_force_long_coeff - brake_force

    # --- Friction circle calculations ---
    # Combined tire forces per axle must stay within the friction circle.
    friction_force_front = self.tire_friction_coeff * weight_front
    friction_force_rear = self.tire_friction_coeff * weight_rear

    input_forces_front = np.hypot(force_front_long, force_front_lat)
    if input_forces_front > friction_force_front and input_forces_front > 0.0:
        scale = friction_force_front / input_forces_front
        force_front_long *= scale
        force_front_lat *= scale

    input_forces_rear = np.hypot(force_rear_long, force_rear_lat)
    if input_forces_rear > friction_force_rear and input_forces_rear > 0.0:
        scale = friction_force_rear / input_forces_rear
        force_rear_long *= scale
        force_rear_lat *= scale

    # --- Calculate Total Forces and Resultant Acceleration ---
    drag_force_lat = - const.DEFAULT_AIR_RESISTANCE_COEFF * self.velocity.x * \
        abs(self.velocity.x) - const.DEFAULT_ROLLING_RESISTANCE_COEFF * \
        self.velocity.x * self.mass
    # Adjust drag based on traction
    drag_force_lat /= (self.tire_friction_coeff)
    total_force_lat = force_front_lat + force_rear_lat + drag_force_lat

    drag_force_long = - const.DEFAULT_AIR_RESISTANCE_COEFF * self.velocity.y * \
        abs(self.velocity.y) - const.DEFAULT_ROLLING_RESISTANCE_COEFF * \
        self.velocity.y * self.mass
    # Adjust drag based on traction
    drag_force_long /= (self.tire_friction_coeff)
    total_force_long = force_front_long + force_rear_long + drag_force_long

    # --- Integrate Physics ---
    # Acceleration = Force / Mass
    accel_lat = total_force_lat / self.mass
    self.velocity.x += accel_lat * dt

    accel_long = total_force_long / self.mass
    self.velocity.y += accel_long * dt

    # Simple anti-overshoot rule: braking while moving forward should
    # stop at zero, not instantly create reverse speed.
    if self.brake > 0.0 and self.prev_velocity.y > 0.0 and self.velocity.y < 0.0:
        self.velocity.y = 0.0

    # Calculate angular velocity based on longitudinal velocity and steering angle
    self.speed = np.hypot(self.velocity.x, self.velocity.y)
    if self.speed > const.MIN_VELOCITY_THRESHOLD:
        speed_factor = 1.0 / \
            (1.0 + const.DEFAULT_UNDERSTEER_GAIN * (self.speed**2))
        self.angular_velocity = (
            const.DEFAULT_STEERING_CORRECTION_COEFF * self.velocity.y / self.wheelbase
        ) * np.tan(np.radians(self.steer_angle)) * speed_factor
    else:
        self.angular_velocity = 0.0

    self.rotation += np.degrees(self.angular_velocity) * dt
    self.rotation = self.rotation % 360  # Keep rotation within [0, 360)

    self.velocity.x *= max(0.0, 1.0 - 4.0 * dt)

    # --- Convert Local Velocity to World Position ---
    rad = np.radians(self.rotation)
    world_vx = np.cos(rad) * self.velocity.x - \
        np.sin(rad) * self.velocity.y
    world_vy = np.sin(rad) * self.velocity.x + \
        np.cos(rad) * self.velocity.y

    self.position.x += world_vx * dt
    self.position.y += world_vy * dt
\end{verbatim}
Model fizyczny bazuje na dynamice pojazdu w układzie lokalnym. Kluczowym elementem jest obliczanie kątów znoszenia opon ($\alpha$), które zależą od prędkości poprzecznej pojazdu, prędkości kątowej oraz geometrii podwozia:
\begin{equation}
    \alpha_{front} = \arctan\left(\frac{v_x + \omega \cdot a}{v_y}\right) - \delta
\end{equation}
gdzie $v$ to składowe prędkości, $\omega$ prędkość kątowa, $a$ odległość osi od środka masy, a $\delta$ kąt skrętu kół.

Siły boczne są limitowane przez tzw. krąg tarcia (ang. \textit{friction circle}), co zostało zaimplementowane poprzez skalowanie wektora sił wypadkowych, jeśli przekraczają one iloczyn współczynnika przyczepności $\mu$ i siły docisku $F_z$:
\begin{equation}
    \sqrt{F_{long}^2 + F_{lat}^2} \leq \mu \cdot F_z
\end{equation}
Dzięki temu model oddaje efekt "płużenia" auta przy zbyt gwałtownym hamowaniu w zakręcie.

\section{Dyskusja osiągniętych wyników}
\subsection{Zalety aplikacji}
\begin{itemize}
    \item \textbf{Modułowość:} Dzięki zastosowaniu ustandaryzowanych protokołów (plik \texttt{actions.py}), logika gry jest oddzielona od implementacji graficznej, co ułatwia testowanie, rozbudowę i potencjalną zmianę biblioteki graficznej bez konieczności modyfikowania logiki gry.
    \item \textbf{Realistyczny model jazdy:} Implementacja uproszczonego \textit{"kręgu tarcia"} sprawia, że samochód zachowuje się w sposób bardziej realistyczny, uwzględniając zarówno siły podłużne, jak i poprzeczne, co przekłada się na bardziej satysfakcjonujące doświadczenie z jazdy, np. auto traci sterowność podczas gwałtownego hamowania przy skręconych kołach.
    \item \textbf{Elastyczność map:} System wczytywania map z plików .json pozwala na łatwe tworzenie nowych torów wyścigowych bez konieczności modyfikowania kodu, a jedynie poprzez edycję danych mapy.
    \item \textbf{Stabilność numeryczna:} Zastosowanie reguły \textit{anti-overshoot} w kodzie fizyki eliminuje drgania samochodu przy zerowych prędkościach, co jest częstym problemem w prostych symulatorach.
\end{itemize}
\subsection{Ograniczenia i potencjalne ulepszenia}
\begin{itemize}
    \item \textbf{Wydajność GUI:} Zastosowanie standardowego obiektu \texttt{tk.Canvas} do renderowania całej sceny może prowadzić do spadków wydajności, zwłaszcza przy większej liczbie obiektów na ekranie. Potencjalnym ulepszeniem byłoby wykorzystanie biblioteki \textit{Pygame} lub \textit{PyOpenGL}, które są zoptymalizowane pod kątem renderowania grafiki 2D i 3D.
    \item \textbf{Brak kolizji:} Obecna implementacja uwzględnia wyłącznie interakcję samochodu z podłożem oraz wypadnięcie z mapy, ale nie implementuje kolizji z potencjalnymi przeszkodami na torze.
    \item \textbf{Brak modelu uszkodzeń:} Samochód nie posiada modelu uszkodzeń, co mogłoby dodać dodatkową warstwę realizmu i strategii do gry, np. poprzez zmniejszenie osiągów pojazdu po kolizji.
\end{itemize}

\end{document}